\chapter{What Is Proved, and What Is Not}\label{ch:ledger}

The book has used four words throughout --- verified, conditional, checked, unverified --- and this
chapter is their ledger. It is also the argument of the book stated once, plainly: a verified
system is not one where everything is proved, but one where the boundary between proved and not is
drawn in the code, reported to the user, and never crossed silently.

\section{The four statuses, once more}

\begin{itemize}[nosep]
  \item \status{verified}: an unconditional theorem, over $\mathbb{R}$ where the rule has a real
    meaning, over $\mathbb{Q}$ for elimination, over the Prop for the order world.
  \item \status{conditional}: a theorem with a side condition whose necessity is itself proved by
    a counterexample. The notebook shows the condition.
  \item \status{checked}: a guess that a later step verifies, and the later step is verified.
  \item \status{unverified}: no theorem, and the note says what would be needed.
\end{itemize}

Appendix~\ref{app:rules} is the table, generated from the same list the engine sends the notebook,
so the book and the software cannot disagree about it.

\section{Termination}

Cell evaluation has no step budget. Every rule of the notebook pipeline decreases the six-tier
ordering on a node whose children are normal (\lc{pipelineOrdered}), for every checker the
pipeline is instantiated with (\lc{pipelineOrderedWith}). Expansion is a structurally recursive
function. Elimination is a loop over columns. The one remaining budget is $\beta$-reduction, whose
termination is undecidable, and it is refused loudly at 1000 steps.

Where the proof needed a property of a subsystem's output it does not see inside, that property is
decided at run time and a failure is an error: a command's output has no command in it; a radical
rule's guard holds; the Kleene chain stabilized. These are boundaries, not budgets. A boundary is a
property of a result; a budget is a count of steps; the difference is whether the failure can be
explained.

\section{Axioms}

Every theorem in \file{engine/} and \file{proofs/} depends on the three standard axioms of Lean's
classical mathematics and no others: propositional extensionality, choice, and quotient soundness.
There are no \lc{sorry}s. \lc{#print axioms} on the top-level theorems is part of the build.

\section{Sizes}

\begin{center}
\begin{tabular}{lr}
  Lean source (engine and proofs) & 11\,389 lines \\
  Theorems and lemmas & 547 \\
  WebAssembly module & 1.8 MB \\
  First reply in the browser & 96 ms \\
\end{tabular}
\end{center}

\section{What is not proved}

The list is in the table, but the pattern is worth stating. What is unproved is either
\emph{vacuous} (a theorem about matrices in a semantics where matrices have no value), an
\emph{abbreviation} (nothing to prove), \emph{symbolic} where the verified path is numeric, or
\emph{metatheory} that the author has not yet written ($\alpha$-equivalence, de Bruijn
commutation, pigeonhole). None of it is a rule that is believed wrong. And every item was on the
open list before it was in this chapter.

\section{What the book argued}

That a learning notebook can show its work and mean it. That the meaning is a theorem per rule,
folded once over the rewriter and instantiated per semantics. That termination without a budget
is an ordering, that the ordering is a design, and that what the ordering cannot afford is better
done as a function than faked as a rule. That a guess plus a verified check is a verified
answer. And that the honest label on a step is worth more than a proof the user cannot see.

The notebook is at the address on the title page. Every ``Try it'' box in this book runs there,
in the browser, on the engine whose theorems this chapter counted.
